%%======================================================================
\section{Background}
\label{sec:background}
%%======================================================================

Before assessing the clinical evidence, three core developments must be
understood: the healthcare workforce crisis that creates demand for robotic
assistance, the humanoid form-factor argument that makes bipedal platforms
specifically relevant, and the LLM/VLA inflection that makes general-purpose
clinical deployment conceivable for the first time. Two supplementary elements
complete the context: a historical trajectory situating current platforms within
the arc of bipedal robotics, and a structured comparison of humanoid against
non-humanoid robotic alternatives.

\subsection{The Healthcare Imperative}
\label{sec:imperative}

The global healthcare workforce is under structural strain that technology alone
cannot resolve, but that technology may help mitigate. The World Health
Organization projects a shortage of approximately 10 million healthcare workers
by 2030, concentrated primarily in low- and middle-income countries but with
significant pressure also on high-income systems where demographic aging compounds
absolute workforce insufficiency~\cite{who2016workforce}. The root cause is well
understood: populations in Organisation for Economic Co-operation and Development
(OECD) countries are aging faster than working-age cohorts are
growing~\cite{oecd2023health}. In Japan (the most extreme case), persons aged 65
and over already represent approximately 30\% of the total population, a
proportion that continues to rise under below-replacement fertility
rates~\cite{japancabinet2023}. The 2023 robotics roadmap from Japan's
Ministry of Economy, Trade and Industry (METI) explicitly identifies
nursing-capable robots as a national strategic target by 2030,
reflecting a governmental judgment that the eldercare workforce gap cannot be
closed by recruitment alone.

Nursing workforce shortages compound this pressure~\cite{lasater2021,aiken2014nurse},
and the care gap operates at two intersecting levels. First, the absolute number
of care hours demanded rises with the prevalence of chronic disease,
multimorbidity, and frailty in an older
population~\cite{barnett2012epidemiology,clegg2013frailty}. A patient with
dementia and comorbid heart failure requires substantially more nursing time
than one with a single, stable
condition~\cite{gopfert2021primary,arling2003cognitive,moth2012chronic}, and the
number of people living with such complex needs is projected to grow
rapidly~\cite{kingston2018projections,prince2015burden,chowdhury2023global}.
Second, the tasks that constitute professional care (physical repositioning of
patients, vital-sign monitoring, medication administration, rehabilitation
exercise guidance, and social engagement) are labor-intensive and resist
substitution by currently deployed information technology. Electronic health
records and remote monitoring shift documentation and surveillance load but do
not remove the need for physical presence at the
bedside~\cite{sinsky2016allocation,sun2020nurses,koehler2018efficacy}.

Automation is best understood as a complement to human care rather than a
replacement for it. The goal is not to displace nurses or physicians but to
extend their reach by delegating tasks that are physically demanding, repetitive,
or logistically routine, freeing clinicians for judgment and human connection.
Three layers of this automation are already established. Wheeled mobile robots
handle medication transport and specimen delivery in
hospitals~\cite{bloss2011mobile,ozkil2009service,fragapane2021autonomous}.
Powered exoskeletons and assistive lifting devices support rehabilitation and
reduce the musculoskeletal strain that patient handling places on
caregivers~\cite{nelson2006evidence,theurel2019occupational,mehrholz2017electromechanical}.
Social and companion robots provide cognitive stimulation and companionship in
dementia care~\cite{pu2019effectiveness,mordoch2013social,wada2007living}. What
remains outside routine clinical deployment is a robot that can perform the
hands-on physical care tasks (patient repositioning, standing assist, equipment
handling, clinical examination) that still require a human body next to the
patient~\cite{sorensen2025humanoid}. It is precisely this layer, which combines
human-scale reach, bimanual manipulation of clinical tools, and
patient-proximate physical assistance, that motivates interest in the humanoid
form factor. Socially assistive robotics~\cite{feilseifer2005defining,riek2017healthcare}
addresses part of the gap through interaction and coaching but not the
physical-care layer, and earlier physical-assistance robots for tasks such as
patient transfer point to the bipedal humanoid as a plausible candidate platform
for that role~\cite{mukai2010riba,ding2014rona,desantis2008atlas}.

\subsection{Why the Humanoid Form Factor}
\label{sec:formfactor}

For the purposes of this survey, a \emph{humanoid robot} is defined as a
bipedal, free-standing robot with a human-proportioned morphology: head, torso,
two arms, and two legs, typically ranging from 150 to 200~cm in height. This
definition explicitly excludes social robots such as Pepper (120~cm, wheeled
base) and NAO (58~cm, child-scale), surgical manipulator arms such as the
da~Vinci system, and wearable exoskeletons. The rationale for this scope
restriction is not hierarchical (social robots and surgical robots are
clinically important and are discussed in Section~\ref{sec:alternatives}), but
conceptual: the form-factor hypothesis concerns whether human-scale bipedal
morphology confers advantages in human-built environments that smaller or
non-mobile systems cannot, specifically adult-height working range, two-handed
reach, compatibility with the spatial scale of beds, cabinets, and doorways, and
the ability to operate clinical tools built for human
hands~\cite{adams2000humanoid,duffy2003anthropomorphism}.

Three distinct lines of argument support the humanoid form factor's suitability
for healthcare deployment.

\textbf{Environmental compatibility.} Hospitals, care homes, and residences are
built for human-sized bipedal operators. Doorways, corridors, stairways, beds,
examination tables, supply cabinets, and clinical instruments are largely
dimensioned for a person roughly 160--180~cm tall with two-handed reach, even
where wheelchairs, stretchers, and carts are also accommodated. In principle, a
bipedal humanoid can therefore navigate these environments without infrastructure
modification: walk through standard doorways, climb stairs, reach across a
standard hospital bed, and operate existing clinical tools (stethoscope,
bag-valve mask, otoscope, ultrasound probe) designed for human
hands~\cite{adams2000humanoid,krotkov2018darpa}.
The
technical demonstration by Atar et al.~\cite{atar2025humanoids} explicitly grounds
this observation: the Unitree G1 demonstrated seven medical procedures using
existing clinical instruments, without facility modification.

\textbf{Neuroscientific basis for human-robot interaction.} Action observation
therapy (AOT), a rehabilitation technique in which patients observe biological
motion to activate sensorimotor circuits, rests on the mirror neuron hypothesis:
observing purposive action activates similar neural circuits as performing that
action~\cite{rizzolatti2004mirror}. Nguyen et al.~\cite{nguyen2024eeg} provide the most direct empirical
evidence available: EEG recordings from healthy participants during observation of
a humanoid robot performing reaching and grasping actions showed sensorimotor
response patterns comparable to those elicited by observing a human agent perform
the same actions. This result is consistent with, though it does not by itself
establish, the stronger hypothesis that a humanoid form contributes materially to
AOT efficacy; comparative HRI studies in which a physically present robot elicits
engagement and task responses that a screen agent does not lend the hypothesis
further support~\cite{bainbridge2011benefits,li2015benefit}.
A complementary behavioral finding comes from Sayis and
Gunes~\cite{sayis2024technology}, who found that the humanoid-robot condition,
but not a voice-only assistant, was associated with improved mood and increased
self-disclosure in a therapeutic writing study (n\,=\,42), a difference they
attribute to the physical embodiment of the humanoid.

\textbf{User acceptance and eye-level interaction.} Patients in clinical settings
are frequently in positions of physical vulnerability and
dependence~\cite{gallagher2008dignity,fuseini2022dignity}.
The SPRING
project~\cite{alamedapineda2024spring} (\emph{scope caveat}: ARI is a wheeled
platform, not bipedal) found that users in a real geriatric day-care facility
explicitly cited upright, eye-level conversational presence as important for
dignified interaction, a form-factor finding directly relevant to bipedal
humanoid design regardless of the mobility substrate. Additional contextual
evidence comes from research on the HRP-4C platform, which demonstrated that
gait naturalness and facial expression generation systematically affected observer
ratings of approachability and social presence~\cite{cao2025humanoid}.

These findings collectively support, but do not yet conclusively prove, the
form-factor hypothesis at clinical scale. The evidence is encouraging; it is not
sufficient to substitute for clinical validation.

\subsection{Why Now: The LLM/VLA Inflection Point}
\label{sec:llm}

Prior to approximately 2022, humanoid robots required explicit, hand-crafted
programming for every behavioral task.
A robot capable of performing a
stethoscope placement procedure had to be programmed, step by step, for that
specific procedure, and deploying the same robot for a second procedure required
a substantially new programming effort~\cite{billard2008robot,schaal1999imitation}.
Learning from demonstration was developed in part to relax this burden, letting a
robot acquire a task from examples rather than explicit
code~\cite{argall2009survey,ravichandar2020recent}, but before the foundation-model
era it still generalized poorly across tasks. This approach fundamentally limited
humanoid deployment in clinical environments, which are intrinsically
unstructured.

The emergence of large language models (LLMs), vision-language models (VLMs), and
vision-language-action (VLA) models has changed this constraint.
GPT-4 and its successors demonstrated that natural language can serve as a
general-purpose task specification interface, capable of interpreting novel
instructions and decomposing them into action plans~\cite{openai2023gpt4}.
Physical Intelligence's $\pi_0$ and $\pi_{0.5}$ models introduced flow-based
VLA architectures mapping language instructions and visual observations directly
to robot actions, trained on thousands of hours of real robot
demonstrations~\cite{black2024pi0}. Google DeepMind's Gemini Robotics applied
multimodal foundation models to general-purpose manipulation with dexterity
approaching human-level performance on standardized
tasks~\cite{geminiroboticsTeam2025}. Universal Manipulation Interface (UMI)
demonstrated that physically deployed robot teaching without robot-specific
teleoperation hardware is viable for generalizable
manipulation~\cite{chi2024umi}. OpenVLA~\cite{kim2024openvla} demonstrated
that open-source, community-accessible VLA models could achieve meaningful
generalization from limited task demonstrations. 3D-VLA~\cite{zhen20243dvla}
extended this to three-dimensional grounding for manipulation in complex
environments. HiRT~\cite{wen2024hirt} demonstrated hierarchical transformers
for whole-body robot control with improved task generalization. RT-2~\cite{brohan2023rt2}
established that vision-language foundation models trained on internet-scale
data transfer task-relevant knowledge to robotic control without task-specific
fine-tuning. Embodied AI surveys covering the healthcare vertical confirm that
this software inflection is accelerating clinical
relevance~\cite{liu2025screens,mir2026embodied,lisondra2026embodied}.

Yuan et al.~\cite{yuan2025integrating} demonstrated this integration pattern in
a dementia care context, combining reinforcement learning with a large language
model to generate adaptive, personalized robot behavior. Although Pepper is a
wheeled robot (not within our bipedal scope), the LLM integration
framework, specifically its ability to model patient-specific
cognitive-emotional states and adapt behavior accordingly, is directly applicable
to bipedal humanoid platforms.

The commercial humanoid wave of 2022--2026 coincides with this LLM inflection,
and the coincidence is not accidental. Twelve of the thirteen commercial humanoid
systems surveyed in Section~\ref{sec:systems} have confirmed or announced
LLM/VLA integration as of 2026. The design assumption of current commercial
humanoids is language-conditioned, generalist behavior: these systems are not
being built as specialized single-task machines, a prerequisite for clinical
task diversity.

\subsection{Historical Evolution}
\label{sec:history}

Figure~\ref{fig:timeline} situates these milestones in the broader thirty-year
arc of humanoid robot development.

\begin{figure*}[!t]
  \centering
  \Description{Dual-track horizontal timeline from 1996 to 2026.
  Upper track (blue): platform and technology milestones including Honda P3
  (1996), ASIMO (2000), DARPA Robotics Challenge (2012), Atlas (2013), Tesla
  Optimus prototype (2022), Unitree G1 release (2024), and thirteen commercial
  platforms competing (2026).
  Lower track (green): healthcare research events including PARO clinical RCT
  (2010), NAO ASD therapy studies (2013), Pepper dementia pilot (2018), SPRING
  project Paris (2022), RHP Friends IREX nursing demo (2023), Atar et al.\ seven
  clinical procedures (2025), and Cho et al.\ humanoid surgical assistant (2026).
  Three shaded eras: locomotion era (1996-2010), manipulation era (2010-2022),
  and commercial scale-up era (2022-present).}
  \includegraphics[width=\textwidth]{figures/tech_timeline.pdf}
  \caption{Technology timeline: humanoid robotics platform milestones (upper
  track, blue) and healthcare research events (lower track, green), 1996--2026.
  Shaded regions demarcate three developmental eras: locomotion proof-of-concept
  (1996--2010), manipulation and integration (2010--2022), and commercial
  scale-up (2022--present). Healthcare research events lag platform capability
  by approximately two years in each era.}
  \label{fig:timeline}
\end{figure*}

The history of humanoid robotics over the past three decades can be organized
into three eras, followed by the nascent healthcare entry period this review
covers.

\textbf{Locomotion proof-of-concept era (1996--2010).} Honda's ASIMO (2000)
demonstrated that dynamic bipedal walking was achievable with then-current
actuation and control technology~\cite{sakagami2002asimo,cao2025humanoid}. The
AIST HRP series (HRP-2 through HRP-5P, 1998--2019) established the research
infrastructure for whole-body motion planning~\cite{kaneko2019hrp5p}. Preview
control of the zero-moment point (ZMP) provided the mathematical foundation for
stable bipedal gait generation~\cite{kajita2003biped}. These systems proved that bipedal locomotion was a
solvable engineering problem, with no meaningful manipulation capability or path
to clinical relevance. ASIMO was retired in 2022.

\textbf{Manipulation and integration era (2010--2020).} The DARPA Robotics
Challenge (DRC, 2012--2015) revealed the gap between locomotion capability and
reliable task execution. Boston Dynamics' hydraulic Atlas became the reference
platform for whole-body manipulation research and produced dynamic capability
demonstrations while revealing that force-controlled manipulation at clinical
precision remained largely unsolved~\cite{cao2025humanoid}. The decade
consolidated learning-based locomotion control, particularly deep reinforcement
learning policies trained in simulation and transferred to
hardware~\cite{duan2022simtoreal,radosavovic2024realworld}. Recent work
demonstrates that transformer-based policies trained on motion-capture data
enable zero-shot locomotion on commercial
humanoids~\cite{radosavovic2024humanoid}. Research platforms such as
TALOS~\cite{stasse2017talos} and iCub3~\cite{dafarra2024icub3} further expanded
the manipulation research base during this period.
Whole-body expressive motion control on commercial humanoids was demonstrated with high fidelity by 2024, enabling
natural interaction behaviors relevant to care
settings~\cite{cheng2024expressive}.

\textbf{Commercial scale-up era (2022--present).} The confluence of electric
actuator maturation, learning-based locomotion policies, and LLM integration
produced the current commercial humanoid wave. Thirteen distinct commercial
humanoid systems reached hardware maturity between 2022 and 2026. Global
investment in humanoid robotics grew rapidly through 2024, driven by the
convergence of electric actuator maturation and LLM integration, with confirmed
commercial deployments concentrated in automotive and logistics
applications~\cite{goldmansachs2024humanoid,ifr2025humanoid}.
As
of March 2026, none of these thirteen systems has been deployed in a confirmed
clinical patient care context.

\textbf{Healthcare entry point (2024--2026).} Atar et al.~\cite{atar2025humanoids}
demonstrated bimanual teleoperation of a Unitree G1 across seven clinical
procedures with approximately 70\% aggregate success, the first systematic
evaluation of a current-generation commercial humanoid against a clinically
meaningful task battery. Cho et al.~\cite{cho2026humanoid} extended this to a
cadaveric endoscopic sinus surgery at Johns Hopkins University, the first
documented use of a modern commercial bipedal humanoid in an actual (cadaveric)
surgical procedure. This survey covers this nascent period.

\subsection{Humanoid Robots vs.\ Other Medical Robots}
\label{sec:alternatives}

A principled comparison with existing clinical robot categories is necessary to
locate the humanoid value proposition precisely.

\subsubsection{Social Robots: NAO, Pepper, PARO}

Social robots~\cite{fong2003survey,dautenhahn2007socially} have the deepest
clinical evidence base in robotics. NAO
(SoftBank/Aldebaran, 58~cm, 25 degrees of freedom (DOF)) has accumulated over
1,000 peer-reviewed
publications, with clinical use spanning autism spectrum disorder therapy,
elderly care, and hospital patient reception~\cite{cabibihan2013robots}. Pepper
has been deployed in approximately 2,000 Japanese healthcare facilities for
dementia cognitive stimulation and companion interaction, with embodied
conversational agent frameworks increasingly integrated for clinical
applications~\cite{provoost2017embodied,sorensen2025humanoid,cunha2025humanoid}.
Observational studies have documented psychosocial benefits of the PARO therapeutic
seal robot in elderly residents~\cite{wada2007living}, and an integrative review
has identified potential for social connectedness in dementia
care~\cite{moyle2017telepresence}, and systematic reviews report benefits for
depression in older adults~\cite{chen2018social}. Long-term deployment in care homes has been
studied for social connectedness and dementia care~\cite{moyle2014dementia}.
These platforms represent the most mature deployments of robots in clinical care
settings to date~\cite{abdi2018scoping}. However, their form-factor constraints are real: NAO at 58~cm cannot
reach adult-height work surfaces or perform physical care tasks for adults; Pepper
on a wheeled base cannot navigate stairways or use standard clinical beds.

\subsubsection{Surgical Robots: da~Vinci}

The da~Vinci Surgical System is FDA-cleared (Class~II, 510(k)) with over 10
million procedures completed across urology, gynecology, and general
surgery~\cite{lanfranco2004robotic,sheetz2020robotic}. Robotic surgery adoption
has grown substantially over the past decade, though adverse event analysis
reveals the importance of systematic safety monitoring even for established
platforms~\cite{alemzadeh2016adverse}. The historical trajectory from laboratory
prototype to clinical instrument provides context for humanoid
deployment~\cite{camarillo2004robotic}. Surgical robots are not competitors to
humanoid clinical assistants; they are complementary systems addressing a
fundamentally different need~\cite{morgan2022robots}. Da~Vinci is stationary, surgeon-controlled, and single-purpose.
A humanoid clinical assistant operating in a post-surgical ward addresses the
nursing and logistics tasks that da~Vinci cannot.

\subsubsection{Exoskeletons: ReWalk, HAL, Ekso Bionics}

Wearable robotic exoskeletons are FDA-cleared (Class~II) for spinal cord injury
rehabilitation~\cite{esquenazi2012rewalk}, with Cochrane-level evidence supporting
electromechanical-assisted training for walking recovery after
stroke~\cite{mehrholz2017electromechanical}.
These devices are established tools in rehabilitation~\cite{sorensen2025humanoid}. Exoskeletons assist the patient's own
movement and are not substitutable with a humanoid assistant: they are worn by
the patient, require near-ambulatory capacity, and cannot lift or reposition a
non-ambulatory patient. The GentleHumanoid framework~\cite{lu2025gentlehumanoid}
demonstrates the humanoid's complementary role: learning upper-body compliance
for sit-to-stand assistance addresses patients for whom exoskeletons are not an
option.

\subsubsection{Wheeled Mobile Robots: TUG, Moxi}

Autonomous wheeled mobile robots address hospital logistics with confirmed
operational deployments at major US health systems~\cite{huang2023intelligent,morgan2022robots}.
These systems are purpose-built for point-to-point transport on flat pre-mapped
surfaces. They cannot navigate stairs, interact with clinical equipment not
specifically engineered for their interfaces, or perform physical patient care
tasks. The humanoid argument relative to wheeled logistics robots is a division
of labor argument: wheeled robots solve the logistics layer effectively;
humanoids are needed for the physical assistance and care layer.

The four alternative categories collectively frame a finding that motivates this
review: no existing clinical robot category combines (a)~general ambulatory
mobility in human-built environments, (b)~bimanual dexterous manipulation of
existing clinical tools, and (c)~natural language-mediated task specification
enabling adaptive behavior in unstructured settings. This combination is the
humanoid value proposition for healthcare. The evidence reviewed in this section
supports the form-factor hypothesis without constituting clinical proof. The
remainder of this survey characterizes where the field stands, what has been
demonstrated, and what must be resolved before that promise can be realized.

